\documentclass[10pt,a4paper]{article}
\usepackage[utf8]{inputenc}\usepackage[T1]{fontenc}\usepackage[brazil]{babel}
\usepackage{geometry,tabularx,booktabs,xcolor,listings}
\usepackage[most]{tcolorbox}
\geometry{margin=1.7cm}\setlength{\parskip}{3pt}\setlength{\parindent}{0pt}
\definecolor{aedblue}{HTML}{17365D}\definecolor{aedorange}{HTML}{D97925}
\lstset{basicstyle=\ttfamily\small,breaklines=true,frame=single,columns=fullflexible}
\newtcolorbox{chave}[1]{colback=blue!3,colframe=aedblue,title={#1},fonttitle=\bfseries,boxsep=1mm,left=2mm,right=2mm,top=1mm,bottom=1mm,before skip=5pt,after skip=5pt}
\newtcolorbox{quadro}[1]{colback=orange!5,colframe=aedorange,title={#1},fonttitle=\bfseries,boxsep=1mm,left=2mm,right=2mm,top=1mm,bottom=1mm,before skip=5pt,after skip=5pt}
\title{AED com LLM --- roteiro docente e chaves dos quatro cenários}
\author{Material do professor --- não distribuir antes das apresentações}
\date{2026.2}
\begin{document}\maketitle

\section*{Ritual simples de cada encontro}
1. Apresente a teoria no quadro. 2. Entregue o cenário e dê 10--15 minutos para o aluno responsável desenhar pergunta, dados, cálculo e resposta. Os outros três atuam como revisores. 3. Compare a proposta do estudante e a resposta sugerida pela IA com a chave abaixo. 4. O responsável implementa em casa uma solução pequena e apresenta o código no encontro seguinte, em até cinco minutos. A chamada não deve interromper a atividade; reserve o tempo de apresentação no início do próximo encontro.

\begin{center}\small
\begin{tabularx}{\textwidth}{p{1.3cm}p{2.6cm}X p{2.4cm}}\toprule
Encontro & Responsável & Cenário & Código apresentado \\\midrule
15/10 & Aluno 1: \rule{1.5cm}{0.4pt} & Cantina: receita e gráfico & 22/10 \\
22/10 & Aluno 2: \rule{1.5cm}{0.4pt} & Loja: métrica e ambiguidade & 29/10 \\
29/10 & Aluno 3: \rule{1.5cm}{0.4pt} & Cafeteria: narrativa com evidência & 05/11 \\
05/11 & Aluno 4: \rule{1.5cm}{0.4pt} & Perguntar, esclarecer ou parar & 19/11 \\\bottomrule
\end{tabularx}\end{center}
12/11 é EXPOFACOL, sem aula regular; por isso o último código é apresentado em 19/11. Não há quinto cenário: o encontro de 19/11 é para a apresentação do código do Aluno 4, integração dos quatro incrementos e avaliação do fluxo agentic.

\begin{quadro}{Esqueleto que vale para os quatro cenários}
Pergunta $\rightarrow$ dado e grão $\rightarrow$ métrica $\rightarrow$ cálculo em Python/SQL $\rightarrow$ validação $\rightarrow$ gráfico/resumo $\rightarrow$ narrativa. A LLM traduz ou explica; o programa calcula e limita as ferramentas. A resposta deve incluir unidade, período e evidência. Se falta definição ou dado, perguntar ou parar.
\end{quadro}

\begin{chave}{Cenário 1 --- Cantina: respostas e código de referência}
Receita por produto: sanduíche $10\times8+12\times8=176$; suco $20\times4+15\times4=140$; fruta $8\times3+10\times3=54$. Maior receita observada: sanduíche, R\$176 nas duas semanas. Isso não prova maior margem, preferência do cliente ou que preço/campanha causou o resultado.

\textbf{Arquitetura no quadro:} CSV (uma linha por produto e semana) $\rightarrow$ validar tipos/ausência $\rightarrow$ calcular receita por linha $\rightarrow$ agrupar por produto $\rightarrow$ barras ordenadas $\rightarrow$ duas frases sustentadas. Sem LLM, o script já automatiza a conta e o gráfico; uma LLM pode interpretar a pergunta e redigir, mas não muda a definição de receita.

\begin{lstlisting}[language=Python]
import pandas as pd
df = pd.DataFrame([
    (1,"Sanduiche",10,8),(1,"Suco",20,4),(1,"Fruta",8,3),
    (2,"Sanduiche",12,8),(2,"Suco",15,4),(2,"Fruta",10,3)
], columns=["semana","produto","unidades","preco"])
df["receita"] = df.unidades * df.preco
totais = df.groupby("produto").receita.sum().sort_values()
print(totais)
totais.plot.bar(title="Receita observada por produto")
\end{lstlisting}
Aceitar \texttt{matplotlib}/\texttt{seaborn} em vez de \texttt{pandas.plot}. Conferir manualmente R\$176, R\$140 e R\$54. Avaliação: grão e fórmula explícitos, resultado correto, gráfico rotulado, conclusão sem causalidade.
\end{chave}

\begin{chave}{Cenário 2 --- Loja: resposta exige métrica}
Faturamento: A = R\$1.000; B = R\$1.200; C = R\$1.050, então B lidera por faturamento. Valor médio por pedido: A $1000/50=\text{R\$}20$; B $1200/80=\text{R\$}15$; C $1050/35=\text{R\$}30$, então C lidera por esse critério. ``Melhor'' não tem resposta única: pedir esclarecimento antes de selecionar.

\textbf{Camada semântica mínima:} \texttt{faturamento = soma do valor recebido na semana}; \texttt{valor\_medio\_pedido = faturamento / número de pedidos}; grão = semana; janela = A--C. Não confundir faturamento com lucro: custos e margem não existem nos dados.

\begin{lstlisting}[language=Python]
dados = {"A": (1000,50), "B": (1200,80), "C": (1050,35)}
ticket = {semana: receita/pedidos
          for semana,(receita,pedidos) in dados.items()}
def responder(metrica):
    if metrica == "faturamento":
        return max(dados, key=lambda s: dados[s][0])
    if metrica == "valor_medio_pedido":
        return max(ticket, key=ticket.get)
    return "Você quer comparar faturamento ou valor médio por pedido?"
\end{lstlisting}
Um mapeamento por regra é suficiente para a demonstração; se houver LLM, ela pode devolver somente uma escolha dentre as duas métricas registradas. O código valida essa escolha. Testes: pergunta por faturamento, por valor médio e ``qual é melhor?''.
\end{chave}

\begin{chave}{Cenário 3 --- Cafeteria: narrativa baseada em evidência}
Receita sobe de R\$1.000 para R\$1.200 ($+20\%$); pedidos sobem de 50 para 80 ($+60\%$); valor médio cai de R\$20 para R\$15 ($-25\%$). Exemplo de narrativa válida: ``Na semana 2, a receita foi 20\% maior, enquanto o valor médio por pedido caiu 25\%. O resumo não informa por que isso aconteceu.'' Não afirmar que uma campanha, produto ou comportamento causou a mudança.

\textbf{Fluxo:} Python calcula indicadores e IDs de evidência (E1 receita; E2 pedidos; E3 valor médio); serializa JSON; LLM redige até três frases; verificador confirma que IDs pertencem a E1--E3; interface mostra narrativa e números originais. Não enviar dados linha a linha quando os agregados bastam.

\begin{lstlisting}[language=Python]
evidencias = {
  "E1": {"metrica":"receita", "s1":1000, "s2":1200,
         "variacao_pct":20},
  "E2": {"metrica":"pedidos", "s1":50, "s2":80,
         "variacao_pct":60},
  "E3": {"metrica":"valor_medio", "s1":20, "s2":15,
         "variacao_pct":-25}}
# Enviar evidencias + instrucao: usar apenas esses fatos;
# devolver claims com evidence_ids. Validar IDs antes de exibir.
import json, os, re
from openai import OpenAI
cliente = OpenAI()  # le OPENAI_API_KEY do ambiente local
resposta = cliente.responses.create(
    model=os.getenv("OPENAI_MODEL", "gpt-5-mini"),
    instructions=("Escreva no maximo tres frases usando so o JSON. "
                  "Nao infira causas. Cite IDs E1, E2 ou E3."),
    input=json.dumps(evidencias, ensure_ascii=False))
texto = resposta.output_text
ids_citados = set(re.findall(r"E\d+", texto))
if not ids_citados.issubset(evidencias):
    raise ValueError("A resposta citou evidencia inexistente")
print(texto)
\end{lstlisting}
O pacote oficial \texttt{openai} precisa estar instalado; a chave vem do ambiente local e o modelo pode ser escolhido pela variável \texttt{OPENAI\_MODEL}. Se não houver conexão, a pessoa pode demonstrar uma resposta simulada, claramente identificada. Testes: números apoiados; pedido de causalidade; ID de evidência inexistente. Nunca publicar a chave no código ou no navegador.
\end{chave}

\begin{chave}{Cenário 4 --- perguntas: responder, esclarecer ou parar}
Com os dados do cenário, A tem resposta: semana B, faturamento R\$1.200. B pede esclarecimento entre faturamento e valor médio; se usar tabela do cenário 2, os líderes são diferentes. C não pode ser respondida causalmente: o conjunto não traz campanha, preço ou outras variáveis explicativas.

\textbf{Ferramentas permitidas:} \texttt{maior\_faturamento()} e \texttt{calcular\_valor\_medio()}. O classificador pode escolher uma função conhecida ou retornar \texttt{clarificar}/\texttt{sem\_evidencia}; nenhuma consulta SQL livre ou ação externa é necessária.

\begin{lstlisting}[language=Python]
def decidir(pergunta):
    q = pergunta.lower()
    if "por que" in q or "causou" in q:
        return "Não há dados para identificar a causa."
    if "melhor" in q and "faturamento" not in q:
        return "Você quer faturamento ou valor médio por pedido?"
    if "faturamento" in q:
        return "Semana B: R$ 1.200 (cálculo sobre a tabela)."
    return "Pergunta não mapeada; não executar ferramenta."
\end{lstlisting}
Este protótipo pode usar regras simples; a discussão agentic é o contrato de ferramenta, a decisão de parar e a validação do resultado, não uma cadeia complexa de agentes. Teste cada classe: suportada, ambígua e sem evidência.
\end{chave}

\begin{quadro}{Rubrica: 0,25 ponto para cada responsável}
Pergunta, grão e métrica corretos: 0,05. Código simples executável: 0,10. Teste e evidência de saída: 0,05. Demonstração de até cinco minutos e resposta a uma pergunta de revisão: 0,05. Soma dos quatro responsáveis: 1,0 ponto. A avaliação não depende de framework, tela elaborada nem plano pago.
\end{quadro}

\textbf{Referências para a conversa no quadro:} Tukey (1977), \emph{Exploratory Data Analysis}; Wickham \& Grolemund (2017), \emph{R for Data Science}; Lewis et al. (2020), Retrieval-Augmented Generation for Knowledge-Intensive NLP Tasks; documentação da Responses API e de ferramentas; documentação de Microsoft Fabric, Tableau Next, Amazon Quick, Looker e Databricks AI/BI. Conceitos avaliados: definição da pergunta, evidência, cálculo reproduzível, resposta sustentada e recusa quando os dados não respondem.
\end{document}
